% Replacement for the current opening "问题描述" subsection and the late
% "动态需求处理" subsection. Insert once near the beginning of Section 2;
% remove the old late dynamic subsection to avoid duplicate labels.

\subsection{问题描述与动态决策过程}

本文研究中国城市群区域—城际配送中的多车场混合车队动态协同问题.
多个车场在保留车辆归属和客户原责任关系的前提下共享配送能力,
车队由燃油车和电动车组成, 每辆实体车在一个运营日内可连续执行多趟任务.
电动车可在共享公共站补能,
也可在首趟出发前或相邻配送趟之间利用车场停留时间按需充电.
电价和电网碳强度按城市、日期和半小时时段给定,
因而同一补电量在不同时空位置会产生不同的电费和间接排放.
模型需要同时确定客户访问顺序、配送趟、实体车衔接、车型、充电位置和充电时刻,
并在载重、硬时间窗、电量、充电桩容量和成员参与条件下
最小化运营成本与碳结算成本之和.

静态计划并不是运营终点.
新增、取消和改量订单到达后,
已经完成、已经发车和正在执行的路径前缀不可撤回;
系统只对尚未执行的客户和车辆后续任务重新优化.
因此, 每次重规划都以“已发生状态+待决策剩余问题”为输入,
而不是从原始算例重新开始.

参照动态需求批处理研究\cite{ref:18,ref:71},
本文采用定时和定量并行的事件触发方式.
设第$m-1$次重规划发生于时刻$\tau_{m-1}$,
其后新到事件的累计需求变动量为$W(\tau_{m-1},u)$,
则第$m$次重规划时刻为
\begin{equation}
\tau_m=\min\left\{
\min\left\{u:W(\tau_{m-1},u)\ge\overline W\right\},
\tau_{m-1}+T^{\mathrm{int}}
\right\}.
\label{eq:trigger}
\end{equation}
累计需求变动达到阈值$\overline W$,
或距上次重规划达到间隔$T^{\mathrm{int}}$时, 触发下一次优化.
在重规划时刻$\tau$, 有效客户集合按已服务、取消和新增事件更新:
\begin{equation}
N^\tau=
\left(N^{\tau-1}\setminus N_{\mathrm{ser}}^\tau
\setminus N_{\mathrm{can}}^\tau\right)
\cup N_{\mathrm{new}}^\tau.
\label{eq:customer_update}
\end{equation}

每辆实体车$k$继承当前位置$o_k^\tau$、最早可用时刻$t_k^\tau$、
剩余载重$\ell_k^\tau$和剩余电量$\varepsilon_k^\tau$.
已完成部分的成本、排放和车场收益按阶段累计:
\begin{equation}
\begin{aligned}
\bar Z^\tau
&=\bar Z^{\tau-1}+Z^{\tau-1},\\
\bar E^\tau
&=\bar E^{\tau-1}+\Delta E^{\tau-1},\\
\bar\Pi_d^\tau
&=\bar\Pi_d^{\tau-1}+\Pi_d^{\mathrm{stage},\tau-1}.
\end{aligned}
\label{eq:accumulation}
\end{equation}
上述状态继承使新方案不能撤回已执行路径、重复使用实体车或重复结算历史排放,
并使后续参与条件建立在成员已实现收益和剩余合作机会之上.

模型作如下基本假设:
1) 每次重规划时, 已到达订单、节点、车辆状态和外生时段参数已知;
未来事件在实际到达前未知.
2) 每个客户由一辆车服务一次, 需求不可拆分.
3) 车辆在配送全程不超过所属车型的最大载重.
4) 车辆行驶速度在构造情景中设为给定值,
能耗采用与车型和载重相关的CMEM系函数\cite{ref:27,ref:36}.
5) 客户具有硬时间窗, 车辆可提前到达后等待, 但不得迟于时间窗上界.
6) 充电功率为电池荷电状态的分段常函数,
每次充电为连续区间并允许部分充电.
7) 电价和电网碳强度为按城市、日期和半小时时段绑定的外生参数;
本文算例为尽量利用公开运营参数的构造情景, 不冒充企业合同观测.
8) 车场收益按执行车辆所属车场归集且不设事后转移支付,
成员是否参与合作由收益底线约束刻画.

模型符号见表\ref{tab:symbols}.
